\begin{helper}[FiberAndDegreeMixedLiftedIntersectionUnionBound]
\label{helper:FiberAndDegreeMixedLiftedIntersectionUnionBound}
Assume
\[
m=\noderef{TwoBiteNaturalReducedVertexCount}(n),\qquad
p=\noderef{TwoBiteEdgeProbability}(1/2,n),\qquad
n\le m^2,\qquad \log n\ge1.
\]
Then the
\(\noderef{TwoBiteEventProbability}(n,m,p)\)-probability of the bad mixed
intersection event
\[
\neg\left(
\forall r,b\in\operatorname{Fin}(m),\
\left|
\noderef{TwoBiteLiftedNeighborhood}(\omega,\operatorname{inl}(r))
\cap
\noderef{TwoBiteLiftedNeighborhood}(\omega,\operatorname{inr}(b))
\right|
\le 100(\log n)^3
\right)
\]
is at most the probability of the bad base-degree event
\[
\neg\left(
\left[\forall r\in\operatorname{Fin}(m),\
\noderef{GraphDegreeCount}(\omega.1,r)\le2pm\right]
\wedge
\left[\forall b\in\operatorname{Fin}(m),\
\noderef{GraphDegreeCount}(\omega.2.1,b)\le2pm\right]
\right)
\]
plus
\[
m^2\exp(-2(\log n)^3).
\]
\end{helper}

\begin{proof}
Let \(A\) be the event that all mixed lifted intersections are at most
\(100(\log n)^3\), and let \(B\) be the event that all red and blue base
degrees are at most \(2pm\), exactly as displayed in the statement.  Put
\[
C=m^2\exp(-2(\log n)^3).
\]
We first prove the initial probability split.  If \(\omega\in A^c\), then
either \(\omega\notin B\), in which case \(\omega\in B^c\), or
\(\omega\in B\), in which case \(\omega\in A^c\cap B\).  Equivalently, for
every sample \(\omega\),
\[
\mathbf 1_{A^c}(\omega)
\le
\mathbf 1_{B^c}(\omega)+\mathbf 1_{A^c\cap B}(\omega).
\tag{1a}
\]
To pass from this pointwise inequality to probabilities, first note that the
specialized value of \(p\) lies in \([0,1]\).  Indeed,
\[
p=\noderef{TwoBiteEdgeProbability}(1/2,n)
=\frac12\sqrt{\frac{\log n}{n}},
\]
and the hypothesis \(\log n\ge1\) implies \(n>0\) and
\(\log n\ge0\).  The elementary inequality \(\log x\le x\) for
\(x>0\), applied to \(x=n\), gives \((\log n)/n\le1\).  Hence
\[
0\le p\le \frac12\le1.
\tag{1b}
\]
By \(\noderef{TwoBiteSampleWeightNonnegative}\), every summand
\(\noderef{TwoBiteSampleWeight}(\omega)\) in the
\(\noderef{TwoBiteEventProbability}\) sums is therefore nonnegative.  Multiply
(1a) by this nonnegative weight and sum over the finite sample space
\(\noderef{TwoBiteSample}(n,m,p)\).  Using the definition of
\(\noderef{TwoBiteEventProbability}\) as the corresponding finite weighted
indicator sum, we obtain
\[
\Pr(A^c)\le \Pr(B^c)+\Pr(A^c\cap B).
\tag{1}
\]

It remains to bound the second term in (1).  Let
\[
E=A^c\cap B.
\]
We verify the hypothesis of
\(\noderef{FiberAndDegreeMixedLiftedIntersectionConditioning}\) for the event
\(E\) with the constant \(C\).  Fix an ordered pair of base graphs
\((G_R,G_B)\), and condition on the cell
\[
\omega.1=G_R,\qquad \omega.2.1=G_B.
\]
If these fixed graphs do not satisfy the degree event \(B\), then no sample
in this conditioning cell can lie in \(E=A^c\cap B\).  Thus the conditional
probability of \(E\) in this cell is \(0\), hence at most \(C\).

Suppose instead that \(G_R,G_B\) satisfy \(B\).  Then, for every \(r,b\), the
two degree bounds required in
\(\noderef{FiberAndDegreeMixedLiftedIntersectionUniformBound}\) hold.  The
specialized hypotheses \(m=\noderef{TwoBiteNaturalReducedVertexCount}(n)\),
\(p=\noderef{TwoBiteEdgeProbability}(1/2,n)\), \(n\le m^2\), and
\(\log n\ge1\) are exactly the additional hypotheses of that node.  Therefore
the conditional probability of \(A^c\) in this cell is bounded by \(C\):
\[
\Pr(A^c\mid G_R,G_B)\le C.
\tag{2}
\]
Since \(E\subseteq A^c\), conditional monotonicity gives
\[
\Pr(E\mid G_R,G_B)\le C
\]
in the good-degree cells as well.  We have proved the same conditional bound
for every pair \((G_R,G_B)\).  Applying
\(\noderef{FiberAndDegreeMixedLiftedIntersectionConditioning}\), which
averages a uniform per-base-graph conditional bound over the base-graph
cells, yields
\[
\Pr(A^c\cap B)=\Pr(E)\le C=m^2\exp(-2(\log n)^3).
\tag{3}
\]
Combining (1) and (3) is exactly the asserted inequality.
\end{proof}
